\section{Appendix}

\subsection{Current implementation details}

The current MGR-SID v1 implementation lives in an isolated experimental branch of the repository and does not alter the default MiniOneRec tokenizer path. The main training and generation entry points are:
\begin{itemize}[leftmargin=1.4em]
  \item \texttt{scripts/experiment\_mgr\_sid\_v1\_train.py}
  \item \texttt{scripts/experiment\_mgr\_sid\_v1\_generate.py}
\end{itemize}

The core implementation is in:
\begin{itemize}[leftmargin=1.4em]
  \item \texttt{src/onerec/experiments/mgr\_sid/train\_v1.py}
  \item \texttt{src/onerec/experiments/mgr\_sid/transplanted\_graph\_bank.py}
  \item \texttt{src/onerec/experiments/mgr\_sid/paper\_transplants.py}
\end{itemize}

\subsection{Graph construction choices in v1}

\begin{table}[h]
  \centering
  \caption{Concrete graph-bank settings used by the current v1 experiments.}
  \label{tab:appendix-graph}
  \begin{tabular}{ll}
    \toprule
    Component & Setting \\
    \midrule
    History window & 10 \\
    Coarse minimum edge weight & 2.0 \\
    Local minimum edge weight & 1.0 \\
    Semantic anchor top-$k$ & 32 \\
    Semantic anchor mixing weight & 0.35 \\
    Spectral rank & 48 \\
    Middle-band eigen ratio & [0.25, 0.65] \\
    Temporal mixing ratio (alt. middle view) & 0.35 \\
    Graph top-$k$ per row in training & 32 \\
    \bottomrule
  \end{tabular}
\end{table}

\subsection{Tokenizer training hyperparameters}

\begin{table}[h]
  \centering
  \caption{Upstream-aligned tokenizer training regime used for the fair comparison in \Cref{tab:main}.}
  \label{tab:appendix-train}
  \begin{tabular}{ll}
    \toprule
    Hyperparameter & Value \\
    \midrule
    Number of quantizers & 3 \\
    Codebook sizes & [256, 256, 256] \\
    Quantized embedding dimension & 32 \\
    Encoder/decoder widths & [2048, 1024, 512, 256, 128, 64] \\
    Optimizer & AdamW \\
    Learning rate & $10^{-3}$ \\
    Weight decay & 0 \\
    Scheduler & constant with warmup \\
    Warmup epochs & 50 \\
    Epochs & 10000 \\
    Batch size & 20480 \\
    Quantization loss weight & 1.0 \\
    $\beta$ & 0.25 \\
    Graph loss weights & $(0.05, 0.15, 0.05)$ \\
    \bottomrule
  \end{tabular}
\end{table}

\subsection{Additional interpretation of the local ambiguity result}

The local ambiguity analysis compares the final baseline index with the final hierarchy-aware index rather than comparing recommendation outputs. This makes the analysis tokenizer-native. The result is therefore best interpreted as evidence about the structure of the tokenizer itself:
\begin{itemize}[leftmargin=1.4em]
  \item the hierarchy-aware tokenizer produces less crowded same-$l_2$ buckets;
  \item the reduction is visible both at catalog level and when weighted by test targets;
  \item the structural change is aligned with the motivating failure mode of local leaf ambiguity.
\end{itemize}

\subsection{Scope of the current draft}

This paper should be read as a tokenizer-stage draft rather than a finished end-to-end paper. The remaining milestones are:
\begin{enumerate}[leftmargin=1.4em]
  \item add repeated-seed verification for the final tokenizer comparison;
  \item evaluate the hierarchy-aware tokenizer on Office under the same upstream-aligned regime;
  \item feed the improved SID into SFT and then RL;
  \item compare against more published tokenizer baselines once downstream metrics are available.
\end{enumerate}
